\chapter{Appendix}

\section*{Project Links}
\textbf{Configured Demo:} \url{https://swiftbite-ryf2.onrender.com} \\
\textbf{GitHub Repository:} \url{https://github.com/abhibitasanya/SwiftBite}

\section{Development Commands}

\subsection*{(a) Frontend}

The frontend is a React, TypeScript, and Vite project located in \texttt{frontend/}. The available commands are:

\begin{verbatim}
cd frontend
npm install
npm run dev
npm run build
\end{verbatim}

\subsection*{(b) Backend}

The Express and TypeScript API is located in \texttt{backend/}. Its package scripts provide development, production build, production start, and lint commands:

\begin{verbatim}
cd backend
npm install
npm run dev
npm run build
npm start
npm run lint
\end{verbatim}

The backend reads database and optional AI configuration from environment variables. The repository includes \texttt{backend/.env.example}; real credentials should be supplied locally and must not be included in the report or source control.

\subsection*{(c) Project Structure}

\begin{verbatim}
SwiftBite/
|-- frontend/
|   |-- src/
|   |   |-- App.tsx
|   |   |-- screens/
|   |   |-- components/
|   |   |-- portals/
|   |   |-- services/
|   |   |-- data/mockData.ts
|   |-- public/
|   |-- package.json
|   |-- vite.config.ts
|
|-- backend/
|   |-- src/
|   |   |-- app.ts
|   |   |-- server.ts
|   |   |-- db.ts
|   |   |-- aiService.ts
|   |   |-- schema.sql
|   |-- package.json
|
|-- Latex/
|-- README.md
\end{verbatim}

\section{Selected Implementation Details}

\subsection*{(a) Frontend Entry Point}

The frontend is mounted with React DOM from \texttt{frontend/src/main.tsx}:

\begin{tcolorbox}[colback=gray!5,colframe=black,title=React Application Entry]
\begin{verbatim}
import React from 'react'
import ReactDOM from 'react-dom/client'
import App from './App'
import './index.css'

ReactDOM.createRoot(document.getElementById('root')!).render(
  <React.StrictMode>
    <App />
  </React.StrictMode>,
)
\end{verbatim}
\end{tcolorbox}

\subsection*{(b) Backend Server Entry Point}

The backend server reads its listening port from the environment, defaulting to port 4000, and creates the Express application:

\begin{tcolorbox}[colback=gray!5,colframe=black,title=Express Server Entry]
\begin{verbatim}
import { createApp } from "./app.js";

const port = Number(process.env.PORT ?? 4000);
const app = createApp();

app.listen(port, () => {
  console.log(`Backend API running on http://localhost:${port}`);
});
\end{verbatim}
\end{tcolorbox}

\subsection*{(c) Registration Validation}

The backend validates that the confirmation value is present and equal to the chosen password. The frontend performs the same comparison before making the request:

\begin{tcolorbox}[colback=gray!5,colframe=black,title=Registration Schema]
\begin{verbatim}
const registerSchema = z.object({
  fullName: z.string().trim().min(2),
  identifier: z.string().min(1),
  password: z.string().min(6),
  confirmPassword: z.string().min(6),
  role: z.enum(["customer", "delivery", "restaurant", "platform"]),
  loginMode: z.enum(["email", "phone"]),
}).refine((data) => data.password === data.confirmPassword, {
  message: "Passwords do not match.",
  path: ["confirmPassword"],
});
\end{verbatim}
\end{tcolorbox}

\subsection*{(d) Database Configuration}

The MySQL connection pool in \texttt{backend/src/db.ts} reads its host, port, username, password, database name, and optional TLS setting from environment variables. The actual password value is intentionally not reproduced here.

\section{Database Queries}

The authoritative schema is \texttt{backend/src/schema.sql}. The following extracts reflect its users and restaurants tables.

\subsection*{(a) Users Table}

\begin{tcolorbox}[colback=gray!5,colframe=black,title=Users Table Schema]
\begin{verbatim}
CREATE TABLE IF NOT EXISTS users (
  id BIGINT UNSIGNED NOT NULL AUTO_INCREMENT,
  full_name VARCHAR(120) NOT NULL,
  identifier VARCHAR(191) NOT NULL,
  role ENUM('customer', 'delivery', 'restaurant', 'platform') NOT NULL,
  password_hash CHAR(64) NOT NULL,
  created_at TIMESTAMP NOT NULL DEFAULT CURRENT_TIMESTAMP,
  updated_at TIMESTAMP NOT NULL DEFAULT CURRENT_TIMESTAMP
    ON UPDATE CURRENT_TIMESTAMP,
  PRIMARY KEY (id),
  UNIQUE KEY uniq_users_identifier_role (identifier, role)
);
\end{verbatim}
\end{tcolorbox}

\subsection*{(b) Restaurants Table}

\begin{tcolorbox}[colback=gray!5,colframe=black,title=Restaurants Table Schema]
\begin{verbatim}
CREATE TABLE IF NOT EXISTS restaurants (
  id BIGINT UNSIGNED NOT NULL AUTO_INCREMENT,
  name VARCHAR(140) NOT NULL,
  cuisine VARCHAR(120) NOT NULL,
  location VARCHAR(160) NOT NULL,
  eta_minutes INT NOT NULL,
  rating DECIMAL(2,1) NOT NULL,
  description VARCHAR(255) NOT NULL,
  featured TINYINT(1) NOT NULL DEFAULT 0,
  PRIMARY KEY (id),
  UNIQUE KEY uniq_restaurants_name (name)
);
\end{verbatim}
\end{tcolorbox}

\subsection*{(c) Login Lookup}

The backend login route uses a parameterized query scoped to the user's identifier and role:

\begin{tcolorbox}[colback=gray!5,colframe=black,title=Login Lookup]
\begin{verbatim}
SELECT id, full_name, identifier, role
FROM users
WHERE LOWER(TRIM(identifier)) = LOWER(TRIM(?))
  AND role = ?
  AND password_hash = SHA2(?, 256)
LIMIT 1
\end{verbatim}
\end{tcolorbox}

\section{Current System Flow}

\begin{itemize}
\item The browser loads the React application built by Vite.
\item The backend is connected to MySQL, and its schema and restaurant API have been verified. The sign-up and sign-in screens send requests to the backend authentication routes; a file-backed authentication fallback remains available when database operations fail.
\item The customer home and menu use bundled sample data, while cart and checkout state are maintained in the frontend.
\item The Place Order action opens a demonstration confirmation screen. Track My Order opens the delivery map, which loads zoomable OpenStreetMap tiles centered on Bhubaneswar; the marker is static and the order flow does not persist an order or process a payment.
\item Chat messages use the local response helper when no API base URL is configured, or the backend chatbot route when configured. Gemini requires a server-side API key; transient primary-model errors receive one retry, a secondary model is attempted after primary failure, and local rule-based replies remain the final fallback.
\item The portal switcher and non-customer dashboards are prototype views and do not provide production role authorization.
\end{itemize}
